% 拉普拉斯-贝尔特拉米算子（综述）
% license CCBYSA3
% type Wiki

本文根据 CC-BY-SA 协议转载翻译自维基百科\href{https://en.wikipedia.org/wiki/Laplace\%E2\%80\%93Beltrami_operator}{相关文章}。

在微分几何中，拉普拉斯–贝尔特拉米算子（Laplace–Beltrami operator）是对拉普拉斯算子（Laplace operator）的一种推广，它作用于定义在欧几里得空间子流形上的函数，更一般地，也可定义在黎曼流形（Riemannian manifold）及伪黎曼流形（pseudo-Riemannian manifold）上。该算子以皮埃尔–西蒙·拉普拉斯（Pierre-Simon Laplace）与欧金尼奥·贝尔特拉米（Eugenio Beltrami）的名字命名。

对于任意定义在欧几里得空间$\mathbb{R}^{n}$上的二次可微实值函数$f$，拉普拉斯算子（又称拉普拉斯算符，Laplacian）将$f$映射为其梯度向量场的散度，即$f$关于$\mathbb{R}^{n}$的一个正交基的$n$个纯二阶偏导数之和。

与拉普拉斯算子类似，拉普拉斯–贝尔特拉米算子也定义为梯度的散度，它是一个线性算子，将函数映射为函数。

该算子可推广至张量场的情形：定义为协变导数（covariant derivative）的散度；或者，也可推广为作用在微分形式上的算子，使用散度与外微分（exterior derivative）定义。所得到的算子称为拉普拉斯–德拉姆算子（Laplace–de Rham operator），以乔治·德拉姆（Georges de Rham）的名字命名。
